\documentclass[11pt]{article}

\usepackage[margin=1in]{geometry}
\usepackage{amsmath}
\usepackage{amssymb}
\usepackage{array}
\usepackage{graphicx}
\usepackage{hyperref}

\title{Boosted Gaussian Witness Kick  -- A Manufactured Solution}

\author{A Adelmann}
\date{\today}

\newcommand{\vect}[1]{\boldsymbol{#1}}
\newcommand{\dd}{\mathrm{d}}
\newcommand{\E}{\vect{E}}
\newcommand{\B}{\vect{B}}
\newcommand{\Pmom}{\vect{P}}
\newcommand{\betav}{\boldsymbol{\beta}}
\newcommand{\gammal}{\gamma}

\begin{document}
\maketitle

\section{Purpose}

This note defines a manufactured solution for a  low-energy witness
e-/e+ pair, kicked by the electromagnetic field of two highly relativistic
and colliding Gaussian electron bunches.  The goal is to guide an OPALX implementation and
establish a order of magnitude estimate.

It is understood that for the detector design, of a $\gamma\text{-}\gamma$ collider, a full self-consistent model is essential. 

Each source bunch is prescribed as a uniformly (non accelerated) moving,
spherical Gaussian charge distribution in its own rest frame. The rest-frame
Coulomb field is known analytically and is Lorentz-transformed to the lab frame
at the witness event. The resulting lab-frame field is then applied to one
witness electron with the same Boris kick convention used by OPALX.

\section{Current OPALX Witness Model}

The present OPALX witness treatment is a source-frame field sampling model.
The source bunch defines the mesh and the field solve.  A witness particle is
not deposited on that mesh; instead its position is translated and rotated into
the source beam-beam frame, the already available mesh field is gathered there,
and the gathered vectors are rotated back before applying the usual particle
kick.

In this approximation the witness is represented only by its spatial coordinate
in the active tracking step.  The relevant mapping is geometric,
\[
  \vect r_s = R_s\left(\vect r_w + (s_w-s_s)\hat{\vect z}\right),
\]
where \(s_s\) and \(s_w\) are the source and witness path lengths and \(R_s\)
is the rigid rotation into the source beam frame.  The sampled fields are then
\[
  \E_{\mathrm{used}} = R_s^T \E_s(\vect r_s),
  \qquad
  \B_{\mathrm{used}} = R_s^T \B_s(\vect r_s).
\]
This is useful as a first implementation because it reuses the \texttt{BeamBeam} mesh,
but it does not treat the field evaluation as a Lorentz-covariant event
calculation.

The full relativistic reference model will eventually starts from the witness event
\[
  x_w^\mu=(c t_w,\vect x_w)
\]
and evaluates the contribution of each source bunch \(j\) at that event.  Let
the source move with almost constant \texttt{BeamBeam} velocity
\[
  \vect v_j = c\betav_j,
  \qquad
  \gamma_j = \frac{1}{\sqrt{1-\betav_j^2}},
\]
and define the event relative to the source centroid by
\[
  \Delta t_j = t_w - t_{j0},
  \qquad
  \Delta \vect x_j = \vect x_w-\vect X_{j0}.
\]
With the decomposition into components parallel and transverse to
\(\betav_j\), the Lorentz transformation into the source rest frame is
\[
  c t'_j =
  \gamma_j\left(c\Delta t_j-\betav_j\cdot\Delta\vect x_j\right),
\]
\[
  \Delta \vect x'_{j,\parallel}
  =
  \gamma_j\left(\Delta \vect x_{j,\parallel}
  - c\betav_j\Delta t_j\right),
  \qquad
  \Delta \vect x'_{j,\perp}=\Delta \vect x_{j,\perp}.
\]
The term \(\betav_j\cdot\Delta\vect x_j/c\) is the relativity-of-simultaneity
piece that is absent in the current geometric sampling model.

For the manufactured Gaussian reference used here, the source rest-frame
charge density is
\[
  \rho'_j(r') =
  \frac{Q_j}{(2\pi)^{3/2}\sigma_j^3}
  \exp\left(-\frac{r'^2}{2\sigma_j^2}\right),
\]
and the electrostatic field is
\[
  \E'_j(\vect r') =
  \frac{Q_j}{4\pi\epsilon_0}
  \frac{\vect r'}{r'^3}
  \left[
    \operatorname{erf}\left(\frac{r'}{\sqrt{2}\sigma_j}\right)
    -
    \sqrt{\frac{2}{\pi}}\frac{r'}{\sigma_j}
    \exp\left(-\frac{r'^2}{2\sigma_j^2}\right)
  \right],
  \qquad
  \B'_j=0.
\]
More generally, the source rest-frame field tensor is transformed back to the
lab/reference frame,
\[
  F^{\mu\nu}_{j,\mathrm{lab}}(x_w)
  =
  \Lambda^\mu{}_{\alpha}(\betav_j)
  F_j^{\prime\,\alpha\beta}(x'_j)
  \Lambda^\nu{}_{\beta}(\betav_j).
\]
In vector form this is
\[
  \E_j =
  \gamma_j(\E'_j-\vect v_j\times\B'_j)
  -
  \frac{\gamma_j^2}{\gamma_j+1}
  \frac{\vect v_j}{c^2}
  (\vect v_j\cdot\E'_j),
\]
\[
  \B_j =
  \gamma_j\left(\B'_j+\frac{\vect v_j\times\E'_j}{c^2}\right)
  -
  \frac{\gamma_j^2}{\gamma_j+1}
  \frac{\vect v_j}{c^2}
  (\vect v_j\cdot\B'_j).
\]
The fields acting on the witness are the superposition
\[
  \E(x_w)=\sum_j \E_j(x_w),
  \qquad
  \B(x_w)=\sum_j \B_j(x_w).
\]
For the OPALX normalized momentum
\(\Pmom=\gamma_w\betav_w\), the corresponding witness equations of motion are
\[
  \frac{\dd \vect x_w}{\dd t}
  =
  c\frac{\Pmom}{\sqrt{1+\Pmom^2}},
  \qquad
  \frac{\dd \Pmom}{\dd t}
  =
  \frac{q}{m c}
  \left[
    \E(x_w)
    +
    c\frac{\Pmom}{\sqrt{1+\Pmom^2}}\times\B(x_w)
  \right].
\]
This is the model the manufactured solution is meant to test against: event
transformation into each source rest frame, field evaluation there, field-tensor
transformation back to the lab/reference frame, and then the witness kick.

During an active \texttt{BEAMBEAM} window, OPALX treats particle container 0 as
the source bunch.  Additional containers listed through
\texttt{WITNESS\_CONTAINERS} are passive: their charge is not deposited onto the
source mesh, and they only sample the field that was already computed from
container 0.

Let the source container path length be \(s_s\), the witness container path
length be \(s_w\), and the witness particle coordinate in its own container frame
be
\[
  \vect r_w = (x_w,y_w,z_w).
\]
The implemented longitudinal mapping into the source frame is
\[
  \Delta s = s_w-s_s,
  \qquad
  z_s = z_w + \Delta s .
\]
With \(R_s\) the rigid rotation from the reference frame into the source beam
frame, the sampled source-frame coordinate is
\[
  \vect r_s
  =
  R_s\left(\vect r_w + \Delta s\,\hat{\vect z}\right).
\]
OPALX then gathers the current source mesh fields at this point:
\[
  \E_s(\vect r_s), \qquad \B_s(\vect r_s).
\]
After gathering, the witness particle positions are restored and the fields are
rotated back with the inverse rigid rotation,
\[
  \E_{\mathrm{used}} = R_s^T \E_s,
  \qquad
  \B_{\mathrm{used}} = R_s^T \B_s.
\]
Those fields are passed to the normal per-container Boris kick.  In OPALX units
the particle momentum is
\[
  \Pmom = \betav_w \gamma_w,
\]
and the continuous Lorentz-force equation represented by the kick is
\[
  \frac{\dd \Pmom}{\dd t}
  =
  \frac{q c}{m c^2}\E
  +
  \frac{q c^2}{m c^2}\betav_w \times \B.
\]

\subsection{Source Field Currently Available in OPALX}

The \texttt{BeamBeam} field calculation uses the binned quasi-static field path.  This
statement includes the case where the adaptive binning has only one current bin:
as long as a \texttt{BINNING} definition is attached to the field solver, OPALX
dispatches to \texttt{computeBinnedSelfFields()} and applies the same
quasi-static Lorentz correction to that single bin.

For a source bin \(b\), OPALX computes the global mean normalized momentum
\[
  \bar{\Pmom}_b = \langle \betav\gamma\rangle_b,
\]
and the bin Lorentz factor and mean velocity
\[
  \gamma_b = \sqrt{1+\bar{\Pmom}_b^2},
  \qquad
  \vect v_b = c\,\frac{\bar{\Pmom}_b}{\gamma_b}.
\]
Before the Poisson solve, the mesh density is transformed to the approximate bin
rest frame by
\[
  \rho'_b = \frac{\rho_b}{\gamma_b},
\]
and the longitudinal mesh spacing is stretched by
\[
  \Delta z' = \gamma_b \Delta z .
\]
After solving the electrostatic problem in this stretched frame, the code
accumulates the lab-frame fields
\[
  \E_b =
  \gamma_b \E'_b
  -
  (\gamma_b-1)
  (\E'_b\cdot \hat{\vect w}_b)\hat{\vect w}_b,
\]
\[
  \B_b =
  \frac{\gamma_b}{c^2}\vect v_b \times \E'_b,
  \qquad
  \hat{\vect w}_b = \frac{\vect v_b}{|\vect v_b|}.
\]
The accumulated fields are then gathered to source particles, and the current
mesh fields can also be sampled by witness containers.

\subsection{What Is Missing}

The witness model is still a source-frame field sampling approximation.  It
rotates and translates witness geometry into the source beam frame, samples the
mesh, and rotates the sampled vectors back.  It does not evaluate the
electromagnetic field tensor at the witness event.

A covariant treatment would start from a witness event
\[
  x^\mu_w = (c t_w,\vect x_w)
\]
and transform that event into the rest frame of each source bunch.  For a source
with velocity \(\vect v_s=c\betav_s\), the parallel part of the event transform is
\[
  t'_s =
  \gamma_s
  \left(
    t_w - \frac{\betav_s\cdot \vect x_w}{c}
  \right),
\]
\[
  \vect x'_{s,\parallel}
  =
  \gamma_s
  \left(
    \vect x_{w,\parallel} - \vect v_s t_w
  \right),
  \qquad
  \vect x'_{s,\perp}=\vect x_{w,\perp}.
\]
The field tensor would then be transformed back to the lab/reference frame,
\[
  F^{\mu\nu}_{\mathrm{lab}}
  =
  \Lambda^\mu{}_\alpha
  F'^{\alpha\beta}_s
  \Lambda^\nu{}_\beta .
\]

In vector notation, for a source rest-frame field \(\E'\), \(\B'\), the lab field
is
\[
  \E =
  \gamma_s(\E' - \vect v_s\times \B')
  -
  \frac{\gamma_s^2}{\gamma_s+1}
  \frac{\vect v_s}{c^2}
  (\vect v_s\cdot \E'),
\]
\[
  \B =
  \gamma_s\left(\B' + \frac{\vect v_s\times \E'}{c^2}\right)
  -
  \frac{\gamma_s^2}{\gamma_s+1}
  \frac{\vect v_s}{c^2}
  (\vect v_s\cdot \B').
\]
The current witness gather does not perform this per-source or per-witness event
transformation.

The missing physics and numerical pieces are:
\begin{itemize}
  \item no Lorentz transformation of the witness event into each source frame;
  \item no relativity-of-simultaneity term in the sampled source coordinate;
  \item no full \(F^{\mu\nu}\) transformation for the sampled field;
  \item no retarded-time or Lienard-Wiechert treatment for accelerated sources;
  \item no deposition of witness charge back onto the source mesh;
  \item no witness-witness or electron-positron mutual space charge;
  \item no separate Lorentz transformation of the already accumulated source
        mesh fields into the low-energy witness frame.
\end{itemize}

\section{Manufactured Boosted-Gaussian Model}

For a controlled implementation target, prescribe each source bunch \(j\) by
\[
  Q_j,\qquad \sigma'_j,\qquad \betav_j,\qquad
  \gamma_j = \frac{1}{\sqrt{1-\betav_j^2}},
\]
and by the lab-frame center worldline
\[
  \vect R_j(t) = \vect R_{j0} + \betav_j c(t-t_0).
\]
At the witness event \((t,\vect r_w)\), define the lab-frame displacement from
the source center,
\[
  \boldsymbol{\ell}_j = \vect r_w - \vect R_j(t).
\]
Split it into parts parallel and perpendicular to the source velocity,
\[
  \boldsymbol{\ell}_{j,\parallel}
  =
  (\boldsymbol{\ell}_j\cdot\hat{\betav}_j)\hat{\betav}_j,
  \qquad
  \boldsymbol{\ell}_{j,\perp}
  =
  \boldsymbol{\ell}_j-\boldsymbol{\ell}_{j,\parallel}.
\]
For a uniformly moving source, the corresponding source-rest coordinate at the
same lab event is
\[
  \vect r'_j
  =
  \boldsymbol{\ell}_{j,\perp}
  +
  \gamma_j \boldsymbol{\ell}_{j,\parallel}.
\]

The rest-frame charge density is a spherical Gaussian,
\[
  \rho'_j(\vect r')
  =
  \frac{Q_j}{(2\pi)^{3/2}{\sigma'_j}^3}
  \exp\left(-\frac{r'^2}{2{\sigma'_j}^2}\right).
\]
The analytic rest-frame electrostatic potential and electric field are
\[
  \phi'_j(\vect r')
  =
  \frac{Q_j}{4\pi\epsilon_0}
  \frac{\operatorname{erf}(u)}{r'},
  \qquad
  u = \frac{r'}{\sqrt{2}\sigma'_j},
\]
\[
  \E'_j(\vect r')
  =
  \frac{Q_j}{4\pi\epsilon_0}
  \frac{
    \operatorname{erf}(u)
    -
    \sqrt{2/\pi}\,(r'/\sigma'_j)
    \exp[-r'^2/(2{\sigma'_j}^2)]
  }{r'^3}
  \vect r',
\]
with the regular limit \(\E'_j(0)=0\).  Since the source is electrostatic in its
own rest frame,
\[
  \B'_j = 0.
\]

Boosting this field into the lab gives the simpler component form
\[
  \E_{j,\parallel} = \E'_{j,\parallel},
  \qquad
  \E_{j,\perp} = \gamma_j \E'_{j,\perp},
\]
\[
  \B_j = \frac{1}{c^2}\vect v_j \times \E_j.
\]
The total manufactured lab field is
\[
  \E = \sum_j \E_j,
  \qquad
  \B = \sum_j \B_j.
\]

\subsection{Witness Kick}

The manufactured witness uses the same normalized momentum convention as OPALX,
\[
  \Pmom = \betav_w\gamma_w .
\]
The Python prototype applies the same Boris kick sequence used in OPALX:

\begin{eqnarray}
 \Pmom^- &=&
  \Pmom^n
  +
  \frac{1}{2}\Delta t\,\frac{q c}{m c^2}\E, \\
\vect t &=&
  \frac{1}{2}\Delta t\,\frac{q c^2}{\gamma^- m c^2}\B,
  \qquad
  \vect s =
  \frac{2\vect t}{1+\vect t\cdot\vect t}, \\
\Pmom' &=& \Pmom^- + \Pmom^- \times \vect t \\
\Pmom^+ &=& \Pmom^- + \Pmom' \times \vect s, \\
\Pmom^{n+1} &=& 
  \Pmom^+
  +
  \frac{1}{2}\Delta t\,\frac{q c}{m c^2}\E .
\end{eqnarray}

\subsection{Initial Manufactured Cases}

The first Python prototype is
\[
  \texttt{sandbox/python/boosted\_gaussian\_witness.py}.
\]
The table below was generated with
\[
\begin{split}
  K_s &= 245~\mathrm{MeV},\\
  \gamma_s &= 480.453039976,\\
  \beta_s &= 0.999997833949,\\
  Q &= -2.002721~\mathrm{nC},\\
  \sigma' &= 1~\mathrm{mm},\\
  \Delta t &= 1~\mathrm{ps},\\
  \Pmom_w^0 &= (0,0,0).
\end{split}
\]
The two symmetric source centers are both at the IP at \(t=0\), moving in
opposite \(z\)-directions.  The witness charge is an electron charge.

\begin{center}
\small
\resizebox{\textwidth}{!}{%
% Generated by boosted_gaussian_witness.py. Do not edit by hand.
\begin{tabular}{lllll}
\hline
case & $\mathbf r_w$ [mm] & $|\mathbf E|$ [GV/m] & $|\mathbf B|$ [T] & $\Delta\mathbf P$ \\
\hline
symmetric center & $(0.000, 0.000, 0.000)$ & $0.000$ & $0.000$ & $(0.000, 0.000, 0.000)$ \\
symmetric off-axis & $(0.500, 0.000, 0.000)$ & $2.135$ & $0.000$ & $(1.253, 0.000, 0.000)$ \\
single boosted source & $(0.500, 0.000, 0.000)$ & $1.067$ & $3.561$ & $(0.575, 0.000, 0.172)$ \\
asymmetric collision & $(0.500, 0.000, 0.000)$ & $1.921$ & $0.712$ & $(1.124, 0.000, 0.061)$ \\
\hline
\end{tabular}

}
\end{center}

The central symmetric case is intentionally a null test:
\[
  \E(0)=0,\qquad \B(0)=0,\qquad \Delta\Pmom=0.
\]
The off-axis symmetric case checks the transverse electric kick while the two
magnetic fields cancel.  The single-source and asymmetric cases give nonzero
magnetic fields and therefore exercise the full \((\E,\B)\) Boris path.

\subsection{Back-to-Back Breit-Wheeler Pair Propagation}

The Python prototype also contains a two-particle witness test intended to be a
minimal model for a low-energy Breit-Wheeler pair born at the interaction point.
The electron and positron are initialized at the same lab event,
\[
  \vect r_{e^-}(0)=\vect r_{e^+}(0)=\vect 0,
\]
with equal kinetic energy
\[
  K_{e^-}=K_{e^+}=313~\mathrm{keV}.
\]
The corresponding normalized momentum magnitude is
\[
  |\Pmom|
  =
  \sqrt{
    \left(1+\frac{K}{m_ec^2}\right)^2-1
  }
  =
  1.26500561229,
  \qquad
  \beta = 0.784487091427.
\]
For the default test direction \(\hat{\vect u}=\hat{\vect x}\),
\[
  \Pmom_{e^-}(0)=|\Pmom|\hat{\vect u},
  \qquad
  \Pmom_{e^+}(0)=-|\Pmom|\hat{\vect u}.
\]

Three trajectories are integrated with identical initial conditions and timestep:
\begin{itemize}
  \item \texttt{symmetric\_collision}: the pair samples the
        Lorentz-transformed fields of two equal boosted Gaussian source bunches;
  \item \texttt{asymmetric\_collision}: the pair samples the same left source
        and a right source with \(80\%\) of the charge;
  \item \texttt{no\_fields}: the same pair is propagated freely with
        \(\E=\B=0\).
\end{itemize}
All cases are written into the same trajectory CSV with a \texttt{scenario}
column.

For the physical-charge illustration
\[
  Q=-2.002721~\mathrm{nC},\qquad
  \Delta t=0.005~\mathrm{ps},\qquad
  T=50~\mathrm{ps},
\]
the exact creation point is a symmetry null:
\[
  \E(0,\vect 0)=\vect 0,
  \qquad
  \B(0,\vect 0)=\vect 0.
\]
In the symmetric collision with the default pure transverse launch, both
particles remain in the \(z=0\) symmetry plane, so the magnetic fields from the
two equal and opposite source bunches cancel along the trajectory.  The
asymmetric collision breaks this cancellation and samples a nonzero magnetic
field even for the same initial pair direction.  The generated kinematic
summary is:

\begin{center}
\small
\resizebox{\textwidth}{!}{%
% Generated by boosted_gaussian_witness.py. Do not edit by hand.
\begin{tabular}{lllllll}
\hline
case & particle & $\mathbf r_f$ [mm] & $\mathbf P_f$ & $K_f$ [keV] & $\Delta\mathbf r$ [$\mu$m] & $\Delta K$ [eV] \\
\hline
symmetric collision & $e^-$ & $(11.759488, 0.000000, 0.000000)$ & $(1.265096, 0.000000, 0.000000)$ & $313.036$ & $(0.323, 0.000, 0.000)$ & $36.248$ \\
symmetric collision & $e^+$ & $(-11.758843, 0.000000, 0.000000)$ & $(-1.264915, 0.000000, 0.000000)$ & $312.964$ & $(0.323, 0.000, 0.000)$ & $-36.246$ \\
asymmetric collision & $e^-$ & $(11.759456, 0.000000, -0.000041)$ & $(1.265087, 0.000000, -0.000004)$ & $313.033$ & $(0.290, 0.000, -0.041)$ & $32.623$ \\
asymmetric collision & $e^+$ & $(-11.758875, 0.000000, 0.000041)$ & $(-1.264924, 0.000000, 0.000004)$ & $312.967$ & $(0.290, 0.000, 0.041)$ & $-32.622$ \\
\hline
\end{tabular}

}
\end{center}

\begin{figure}[!htp]
  \centering
  \includegraphics[width=\textwidth]{../data/boosted_gaussian_witness_physical_field_t0.png}
  \caption{Magnitude of the manufactured boosted-Gaussian electric field at
  pair creation.  The white marker is the interaction point.}
\end{figure}

\begin{figure}[!htp]
  \centering
  \includegraphics[width=\textwidth]{../data/boosted_gaussian_witness_physical_paths.png}
  \caption{Back-to-back \(e^-/e^+\) propagation with and without the colliding
  source fields.  The lower-right panel shows the field-on trajectory minus the
  free trajectory.}
\end{figure}

\section{Implementation Use in OPALX}

The recommended implementation sequence is:
\begin{enumerate}
  \item Keep the Python evaluator as the reference for exact one-witness values.
  \item Add an OPALX unit or regression test with one prescribed source and one
        witness particle.
  \item First compare the field sampled by OPALX against the Python
        \((\E,\B)\) value at the same event.
  \item Then compare the witness momentum after one kick against the Python
        Boris update.
  \item Only after this single-particle manufactured test passes should the
        implementation be generalized to many witnesses and mesh interpolation.
\end{enumerate}

\end{document}
